%% ACM SIGSPATIAL 2026 Demo Paper
%% Format: sigconf, two-column, max 4 pages including references
%% Single-blind review: author names and affiliations included
\documentclass[sigconf]{acmart}

\usepackage{enumitem}
\usepackage{booktabs}

\setcopyright{none}
\acmConference[SIGSPATIAL '26]{34th ACM SIGSPATIAL International Conference
  on Advances in Geographic Information Systems}{November 2026}{Atlanta, GA, USA}
\acmYear{2026}
\acmDOI{}
\acmISBN{}

\begin{CCSXML}
<ccs2012>
   <concept>
       <concept_id>10002951.10003227.10003236</concept_id>
       <concept_desc>Information systems~Information extraction</concept_desc>
       <concept_significance>500</concept_significance>
   </concept>
   <concept>
       <concept_id>10010405.10010469.10010474</concept_id>
       <concept_desc>Applied computing~Environmental sciences</concept_desc>
       <concept_significance>300</concept_significance>
   </concept>
   <concept>
       <concept_id>10002951.10003227.10003351</concept_id>
       <concept_desc>Information systems~Document structure</concept_desc>
       <concept_significance>100</concept_significance>
   </concept>
</ccs2012>
\end{CCSXML}

\ccsdesc[500]{Information systems~Information extraction}
\ccsdesc[300]{Applied computing~Environmental sciences}
\ccsdesc[100]{Information systems~Document structure}

%% -----------------------------------------------------------------------
\begin{document}

\title{PolarKD Dataset Harvester: Automated Extraction of Geospatial Dataset
  References from Polar Science Literature [Demo]}

\author{Gowtham Vuppaladhadiam}
\affiliation{%
  \institution{University of North Texas}
  \city{Denton}
  \state{TX}
  \country{USA}
}
\email{gowthamvuppaladhadiam@my.unt.edu}

\author{Ajith Kumar Dugyala}
\affiliation{%
  \institution{University of North Texas}
  \city{Denton}
  \state{TX}
  \country{USA}
}
\email{AjithKumarDugyala@my.unt.edu}

\author{Sahara Ali}
\affiliation{%
  \institution{University of North Texas}
  \city{Denton}
  \state{TX}
  \country{USA}
}
\email{sahara.ali@unt.edu}

%% -----------------------------------------------------------------------
\begin{abstract}
Reanalysis products, satellite records, and in-situ observations are cited in
polar science papers in many ways: sometimes with a DOI, sometimes with a
repository accession number, and often just by name with no identifier at all.
Finding all of them automatically is harder than it sounds. We built the PolarKD
Dataset Harvester, a two-pass extraction pipeline that pulls structured dataset
references out of polar science PDFs. The first pass uses deterministic regular
expressions to catch six types of explicit identifiers: URLs, DOIs, PANGAEA
accessions, Zenodo records, NSIDC accessions, and named satellite and reanalysis
products. The second pass sends text chunks to a language model to catch datasets
mentioned by name only, using a signal-based pre-filter to avoid unnecessary API
calls. The system runs on Azure OpenAI, OpenAI, or a local Ollama instance, and
falls back automatically when a backend is unavailable. At the demo, attendees
upload their own papers and watch the extraction run live inside the PolarKD
knowledge graph platform.
\end{abstract}

\keywords{dataset extraction, polar science, geospatial data, scientific literature
  mining, knowledge graph, PDF parsing}

\maketitle

%% -----------------------------------------------------------------------
\section{Introduction}

Reproducibility in polar science depends on knowing which datasets a paper actually
used. A single study might draw on ERA5 reanalysis, MODIS sea ice concentration,
Argo float profiles, and CMIP6 model output at the same time. Tracking those
dependencies manually across hundreds of papers is not practical. Building a tool
that does it automatically is harder than it first appears.

The core problem is that the same dataset can show up in many different forms. ERA5,
for example, might appear as the product name ``ERA5,'' as the DOI
10.24381/cds.adbb2d47, as a URL to the Copernicus Climate Data Store, or buried
inside a Data Availability statement with no identifier at all \cite{era5}. An
extraction tool needs to handle all of those cases without generating too many false
positives. Instrument model names, journal DOIs, figure cross-references, and funding
acknowledgements all appear in similar contexts and cause constant noise.

Polar science also has its own data ecosystem. Repositories like PANGAEA \cite{pangaea},
NSIDC, Zenodo \cite{zenodo}, the Arctic Data Center, and the Copernicus Marine
Service each have their own DOI prefix conventions and accession number formats.
General-purpose scientific NLP tools \cite{luan2018,wadden2019,scibert} are not
tuned to these patterns, and they show it.

We built the PolarKD Dataset Harvester as a module inside the larger PolarKD
(Polar Knowledge Discovery) system. Given any polar science PDF, it returns a
structured list of \texttt{DatasetRef} objects containing the dataset name, URL,
DOI, accession number, inferred repository, the sentence it came from, and a flag
indicating whether this is the paper's single central dataset or a secondary
reference. Most ten-page papers run through in under 30 seconds.

%% -----------------------------------------------------------------------
\section{System Overview}

PolarKD is a Streamlit application we run on a shared research server at the
University of North Texas. It handles keyword extraction (TF-IDF, YAKE, KeyBERT),
semantic relation extraction via local Ollama LLMs, Neo4j knowledge graph
construction, and document Q\&A with RAG. The Dataset Harvester sits inside it as a
self-contained submodule (\texttt{dataset\_harvester/extractor.py}) but can also be
called directly:

\begin{verbatim}
refs = extract_from_pdf("paper.pdf")
\end{verbatim}

\noindent Each returned \texttt{DatasetRef} carries:

\begin{itemize}[noitemsep, topsep=2pt]
  \item \textbf{name}: product name (e.g., ERA5, PANGAEA.940320)
  \item \textbf{url}: direct download or landing URL, if found
  \item \textbf{doi}: data DOI string, if found
  \item \textbf{accession}: repository-specific accession number
  \item \textbf{repository\_hint}: inferred repository (e.g., pangaea, nsidc)
  \item \textbf{raw\_citation}: the sentence in the paper where it appeared
  \item \textbf{is\_primary}: \texttt{True} only for the single dataset the paper
    is centrally about — what it introduces or presents as its main record
\end{itemize}

Both passes run on the same extracted text and feed into a shared deduplication
step before output. The pipeline is shown in Figure~\ref{fig:pipeline}.

\begin{figure}[t]
\centering
\small
\renewcommand{\arraystretch}{1.2}
\begin{tabular}{|c|}
\hline
\\[-1.5ex]
\textbf{PDF Input} \\[3pt]
$\downarrow$ \\[2pt]
\textbf{Text Extraction} \\
{\footnotesize pdfplumber $\cdot$ two-column detection $\cdot$ URL reconstruction} \\[3pt]
$\downarrow$ \\[2pt]
\textbf{Pass 1: Regex} \\
{\footnotesize URLs $\cdot$ DOIs $\cdot$ Accessions $\cdot$ Named products} \\[3pt]
$\downarrow$ \\[2pt]
\textbf{Pass 2: LLM} \\
{\footnotesize Truncate $\cdot$ chunk $\cdot$ signal filter $\cdot$ call $\cdot$ parse $\cdot$ validate} \\[3pt]
$\downarrow$ \\[2pt]
\textbf{Merge, Deduplicate, Output} \\
{\footnotesize DatasetRef list (name, url, doi, accession, is\_primary)} \\[2pt]
\hline
\end{tabular}
\caption{PolarKD Dataset Harvester pipeline.}
\label{fig:pipeline}
\end{figure}

%% -----------------------------------------------------------------------
\section{Technical Approach}

\subsection{PDF Parsing}

We use \texttt{pdfplumber} \cite{pdfplumber} to extract text page by page. For
two-column papers, which are the norm in polar science journals, we split each page
at its horizontal midpoint and process the left and right columns separately before
joining them. Without this step, text from the two columns gets interleaved, which
breaks both the regex patterns and the LLM chunking. We also stitch back URLs that
got split across a line break during PDF rendering. If the total extracted text is
under 3,000 characters, the document is flagged as likely scanned with no embedded
text layer.

\subsection{Pass 1: Deterministic Regex Extraction}

The regex pass scans the full document for six types of identifiers:

\begin{enumerate}[noitemsep, topsep=2pt]
  \item \textbf{URLs}: kept only if the domain is in our list of known data
    repositories, or if the URL appears within 150 characters of language like
    ``available at'' or ``downloaded from.'' Journal publisher domains are excluded.
  \item \textbf{DOIs}: the prefix is checked against a whitelist of 14 data
    repository DOI prefixes (Table~\ref{tab:doi}). Any DOI with a journal publisher
    prefix is discarded.
  \item \textbf{PANGAEA accession numbers}: matched by the pattern
    \texttt{PANGAEA.\{5--7 digits\}}.
  \item \textbf{Zenodo record IDs}: matched by
    \texttt{zenodo.org/record/\{digits\}} or \texttt{zenodo.org/records/\{digits\}}.
  \item \textbf{NSIDC accession numbers}: matched by \texttt{NSIDC-\{4 digits\}}.
  \item \textbf{Named products}: a compiled regex over a list of more than 100
    polar science product names, matched case-insensitively. The list covers
    reanalysis (ERA5, JRA-55, MERRA-2, NCEP/NCAR, CFSR), satellite missions
    (MODIS, VIIRS, AMSR2, ICESat-2, CryoSat-2, Sentinel-1 through 6, GRACE,
    GRACE-FO, CALIPSO, SMAP, GEDI), ocean reanalysis (GLORYS, OSTIA, ORAS5,
    ECCO, HYCOM), precipitation datasets (GPM, IMERG, GPCP, TRMM, CMORPH),
    climate model ensembles (CMIP3, CMIP5, CMIP6, MPI-ESM, CESM, CESM2, E3SM),
    and observation networks (FLUXNET, AmeriFlux, ICOS, SOCAT, GLODAP, ARGO),
    among others.
\end{enumerate}

For every match, we also inspect the surrounding 180 characters for phrases like
``we used'' or ``data were downloaded from,'' and check whether the match falls
inside a Data Availability section. Either condition sets the \texttt{used\_in\_study}
flag on that reference.

\begin{table}[t]
  \caption{Data-repository DOI prefixes retained by the regex pass.}
  \label{tab:doi}
  \small
  \setlength{\tabcolsep}{4pt}
  \begin{tabular}{ll}
    \toprule
    \textbf{Prefix} & \textbf{Repository} \\
    \midrule
    10.1594  & PANGAEA \\
    10.5281  & Zenodo \\
    10.7289  & NOAA NCEI \\
    10.5067  & NASA EOSDIS / NSIDC \\
    10.18739 & Arctic Data Center \\
    10.6073  & BCO-DMO \\
    10.6084  & Figshare \\
    10.5061  & Dryad \\
    10.3334  & ORNL DAAC \\
    10.5285  & BODC \\
    10.48670 & Copernicus Marine Service \\
    10.24381 & Copernicus Climate Data Store \\
    10.5439  & ARM Facility \\
    10.22008 & GEUS / Greenland Data Portals \\
    \bottomrule
  \end{tabular}
\end{table}

\subsection{Pass 2: LLM-Based Extraction}

The second pass finds datasets that are mentioned only by name, where a regex
cannot help.

\textbf{Pre-truncation.}
We cut the text at the first occurrence of ``References,'' ``Bibliography,''
``Acknowledgements,'' ``Acknowledgments,'' ``Supplementary,'' ``Funding,'' or
``Appendix.'' The cut is only applied if that marker appears past 20\% of the
document length (at least 20,000 characters). The guard prevents a
``Funding'' keyword in a page header from accidentally removing the entire paper
body.

\textbf{Chunking.}
The truncated text is split into 3,000-character chunks with 400-character overlap,
so dataset mentions near a chunk boundary are not missed.

\textbf{Signal filter.}
Before calling the LLM, each chunk is tested against a signal regex. If there is
no match, the chunk is skipped entirely. Signals include data language
(``dataset,'' ``accession,'' ``downloaded from,'' ``data availability,''
``code and data''), repository names (``zenodo,'' ``pangaea,'' ``figshare,''
``dryad,'' ``dataverse''), the full named-product list from Pass 1 plus additional
terms (CloudSat, WRF, NEMO, ROMS, FESOM, GlobSnow, and others), observation
descriptors (``reanalysis,'' ``in situ,'' ``buoy,'' ``CTD,'' ``mooring,''
``ice core,'' ``eddy covariance''), field terms (``permafrost,'' ``tundra,''
``peatland,'' ``catchment,'' ``flux tower''), and climate model identifiers (RCP,
SSP, GCM, HadGEM, CESM, MPI-ESM, GFDL, and others). Skipping quiet chunks
substantially cuts API cost.

\textbf{LLM call.}
Each signal-passing chunk is sent to the LLM with a prompt asking for a JSON array
of dataset objects with the fields: name, url, doi, accession, repository\_hint,
raw\_citation, and is\_primary. Temperature is set to 0.

\textbf{Four-layer JSON parser.}
LLM output is not always clean JSON. We try four strategies in sequence:
(1) strip markdown fences and call \texttt{json.loads()};
(2) pull out the outermost \texttt{[...]} block with a regex and parse that;
(3) find individual \texttt{\{...\}} objects if no array brackets exist;
(4) if nothing works, log the chunk as a JSON error and move on.
Layer 4 was added specifically to handle models like llama3 that wrap their JSON
output in explanatory prose.

\textbf{False-positive filter.}
Every name the LLM returns is checked before being accepted. We reject names that
match: lab instrument terms (analyzer, spectrometer, chromatograph, centrifuge);
brand names (Whatman, Sigma-Aldrich, Thermo Fisher, Shimadzu, Agilent, Bruker);
journal names (Biogeosciences, Geophysical Research Letters, Nature Climate Change,
Science Advances); figure and table references (``Fig.~3a,'' ``Table~2,''
``Supplementary Fig.''); author citations containing ``et al.''; funding body names
(NSF, NSERC, NIH, ERC, DFG) and grant strings; acknowledgement strings longer than
250 characters; and bare measurement terms like DOC, DIC, or ``samples.'' Names
under three characters are also rejected.

\subsection{Backend Selection and Fallback}

At startup, the extractor checks for an LLM backend in this order: Azure OpenAI
(gpt-4.1-mini by default, configurable via environment variables), then
OpenAI (gpt-4o-mini), then local Ollama. For Ollama, it queries the
\texttt{/api/tags} endpoint and picks the first available model from the list:
mistral:latest, mistral:7b, llama3:latest, qwen2.5:7b, gemma3:12b. If no backend
is found at all, the LLM pass is skipped and we return only the regex results.
Errors in individual chunks — JSON parse failures, HTTP timeouts — are logged and
skipped without stopping the overall run.

%% -----------------------------------------------------------------------
\section{Demonstration}

The demo runs on our shared UNT research server, reached by SSH tunnel. Nothing
needs to be installed on the attendee side.

We give attendees a set of pre-loaded polar science PDFs spanning cryosphere
observations, physical oceanography, and climate modelling. They can also upload
their own paper. After clicking ``Extract Datasets,'' the Streamlit interface shows:

\begin{enumerate}[noitemsep, topsep=2pt]
  \item \textbf{Extraction progress}: chunks processed, signal-pass rate, LLM call
    count.
  \item \textbf{Pass 1 output}: each regex-found entry labeled by type (URL, DOI,
    PANGAEA accession, named product) with its verbatim source sentence.
  \item \textbf{Pass 2 output}: LLM-recovered name-only datasets with source
    sentences.
  \item \textbf{Merged DatasetRef list}: deduplicated, classified as primary or
    cited, with repository hints and download status (direct download, collection,
    login required, or homepage only).
  \item \textbf{Collection expansion}: for PANGAEA parent records, the system lists
    all child datasets and checks each for download availability via parallel HTTP
    HEAD requests.
\end{enumerate}

Attendees can then open the Knowledge Graph view, where primary datasets appear as
green nodes and cited datasets as blue nodes, connected to the keyword graph the LLM
relation module builds in parallel.

We use three contrasting papers to walk through the system. First, a fieldwork paper
that names more than 15 satellite products with no DOIs or URLs: this puts the LLM
pass and signal filter to work. Second, a data paper deposited to PANGAEA with
explicit DOIs and accession numbers: this exercises the regex pass and PANGAEA
collection expansion. Third, a modelling study referencing CMIP6 outputs and ERA5
forcing: this is where the named-product matching and false-positive filter both get
tested.

%% -----------------------------------------------------------------------
\section{Related Work}

There has been substantial work on extracting structured information from scientific
text, including named entities \cite{luan2018}, relations and events \cite{wadden2019},
and domain-adapted language models \cite{scibert}. Dataset-specific extraction has
received less attention. Google Dataset Search \cite{datasetsearch} takes a different
approach, relying on schema.org metadata embedded in web pages rather than mining
PDF text. The CORD-19 project \cite{cord19} showed what large-scale scientific NLP
can do, but it targeted biomedical literature and did not deal with the
repository-specific identifier conventions that matter in geoscience.

What sets this harvester apart is the combination of a fast deterministic pass for
identifier-rich text with an LLM pass for name-only mentions, together with polar
science specificity: the DOI prefix whitelist, the named-product vocabulary, and the
false-positive rules are all tuned to polar data repositories. Because Ollama is a
supported backend, the pipeline can also run on an air-gapped or restricted research
server without any cloud API access.

%% -----------------------------------------------------------------------
\section{Conclusion}

The PolarKD Dataset Harvester pulls geospatial dataset references from polar science
PDFs using a regex pass for explicit identifiers and an LLM pass for name-only
mentions. The signal-based chunk filter reduces API cost, the four-layer JSON parser
handles messy LLM output, and the false-positive filter keeps noise low. The system
degrades gracefully when cloud APIs are unavailable. We plan to extend the
named-product vocabulary to cover more geoscience domains beyond polar research, and
to build a cross-paper co-citation graph that shows which datasets the polar science
community relies on most.

%% -----------------------------------------------------------------------
\bibliographystyle{ACM-Reference-Format}
\bibliography{references}

\end{document}
